\section*{Spivak Capítulo 3, Problema 16}
\addcontentsline{toc}{section}{Spivak Capítulo 3, Problema 16}

Se supone que $f:\mathbb{R}\to\mathbb{R}$ satisface
\[ f(x+y)=f(x)+f(y) \qquad\text{para todos }x,y. \tag{*}
\]

\textbf{(a)} Para demostrar que
\[ f(x_1+\cdots+x_n)=f(x_1)+\cdots+f(x_n) \]
basta una inducción sobre $n$. El caso $n=1$ es trivial. Si la propiedad vale
para $n$, entonces aplicando (*) con $x=x_1+\cdots+x_n$ e $y=x_{n+1}$ obtenemos
\begin{align*}
f((x_1+\cdots+x_n)+x_{n+1}) &= f(x_1+\cdots+x_n)+f(x_{n+1}) \\
&= (f(x_1)+\cdots+f(x_n))+f(x_{n+1}),
\end{align*}
que es la fórmula para $n+1$. Concluimos por inducción.

\textbf{(b)} Sea $c=f(1)$. En la indicación se sugiere pensar primero en cómo
debe ser $c$; de hecho, tomamos $x=1$ en (*) y obtenemos
$$f(1+y)=f(1)+f(y)=c+f(y).$$
Sustituyendo $y=0$ en (*) da $f(0)=0$ y de allí $c=f(1)$ es coherente.

Para $m$ entero positivo repetimos la aditividad:
$$f(m)=f(1+\cdots+1)=mf(1)=mc.$$  Aquí empleamos la parte (a) con todos
los sumandos iguales a $1$. Si $m$ es entero negativo escribimos $m=-k$ y
entonces
$$0=f(0)=f(m+(-m))=f(m)+f(-m),$$
de donde $f(-m)=-f(m)=-mc$.

Para la inversa de un entero, si $m\neq0$ entonces
$$0=f(m+\frac1m)=f(m)+f(\tfrac1m)=mc+f(\tfrac1m),$$
lo que da $f(1/m)=-mc = c\cdot(1/m)$.

Finalmente, para cualquier $x=p/q$ racional (con $q>0$) usamos
$$f(x)=f\!igl(\tfrac{1}{q}+\cdots+\tfrac{1}{q}\bigr)\text{ ($p$ sumandos)}
= p f(1/q) = p(c/q)= cx.
$$
Si $q$ fuera negativo absorvemos el signo en $p$ para reducir al caso
anterior. Esta demostración muestra que para todo racional $r$ se tiene
$$f(r)=cr.$$

No se exige demostrar la igualdad para irracionales; sólo se necesitaba
establecerla en los racionales. Nótese que la constante $c$ queda única y
está determinada por $f(1)$. 

\textbf{Ejemplo.} Consideremos la función $f(x)=3x$. Es clara que
$$f(x+y)=3(x+y)=3x+3y=f(x)+f(y),$$
por lo que cumple la hipótesis. Aquí $c=f(1)=3$ y la conclusión de (b)
se verifica enseguida: para cualquier racional $r$ tenemos
$f(r)=3r = c r$. Si se desea, la parte (a) se observa también: la suma
de varios argumentos da
$$f(x_1+\cdots+x_n)=3(x_1+\cdots+x_n)=3x_1+\cdots+3x_n
= f(x_1)+\cdots+f(x_n).$$
Este ejemplo muestra que la teoría es compatible con la intuición de
una función lineal con pendiente $c$.

Con esto se completa la parte (b) del problema.